\section{Robustness}
\label{sec:robustness}

The identifying assumption is that, conditional on bank and quarter
fixed effects, the size of the disclosed Day-One adjustment is
unrelated to whatever else was moving deposits in 2023.  Two
cross-sectional facts about the treatment support that assumption, and
this section documents them.  The first is general: the adjustment is
largely unexplained by anything observable about the bank before the
disclosure, so high-CECL banks are not simply the observably weak
banks relabeled.  The second is specific to the setting: because
community banks adopted CECL in the quarter of the Silicon Valley Bank
failure, the sharpest possible confound is that the adjustment might
proxy for exposure to the 2023 stress, and depositors were reacting to
that exposure rather than to the disclosure.  I show the adjustment is
essentially orthogonal to the balance-sheet characteristics that
governed stress exposure, so the stress cannot manufacture the
cross-sectional pattern the paper documents.

\subsection{What Predicts the Day-One Adjustment?}
\label{sec:robustness:predictors}

Table~\ref{tab:cecl_predictors_cross_section} regresses
\texttt{CECL Adj} on 2022Q4 bank characteristics: balance-sheet and
credit-quality fundamentals in column~(1), deposit composition and
eight-quarter performance trends in column~(2), and local
deposit-market concentration in column~(3).  The adjustment loads on
NPL ratios, size, capital, and deteriorating profitability trends in
the expected directions, consistent with more credit-sensitive
portfolios requiring larger expected-loss reserves.  Among funding
variables, the time-deposit and brokered shares enter positively but
with small coefficients.  Adjusted $R^2$ reaches roughly 1.4 percent
in the most saturated specification.  Low explanatory power does not
prove exogeneity, but it establishes that high-CECL banks are not
simply the observably weak or deposit-fragile banks relabeled: almost
all of the variation in the disclosed charge is orthogonal to
everything a depositor could have observed before the disclosure.

\subsection{The Adjustment Is Not an SVB-Exposure Proxy}
\label{sec:robustness:svb}

The remaining concern specific to 2023 is that the adjustment might
proxy for vulnerability to the March 2023 stress.  If it did, banks
disclosing larger charges would also be the banks carrying the
balance-sheet exposures that identified run-prone institutions in 2023
\citep{jiang2024monetary}.  Figure~\ref{fig:svb_exposure} plots three
such exposures, measured at 2022Q4, across deciles of the disclosed
adjustment: unrealized securities losses (Panel~A), the weighted
average maturity of loans and debt securities (Panel~B), and the
composite run-risk measure discussed below (Panel~C).  All three
profiles are flat.  A joint test that the decile means are equal is
not rejected for any of them ($p=0.77$, $0.27$, and $0.42$
respectively), and the decile index explains under one percent of the
cross-sectional variation in each case.  The deciles are formed among
the 1,383 banks with a strictly positive adjustment, the variation the
continuous treatment uses; 53~percent of the full sample discloses an
adjustment of exactly zero, which would otherwise make the lower
deciles arbitrary.

Panel~C uses the composite run-risk measure of
\citet{huberdeau2025adverse}, who summarize a bank's vulnerability as
the expected loss of its uninsured depositors: promised debt over
marked-to-market assets, multiplied by the runnable share of that
debt.  I approximate their Balance Sheet Risk at 2022Q4 as uninsured
deposits plus short-term borrowings over total assets net of
unrealized securities losses.  Conditioning does not help the stress
story either.  In cross-sectional regressions of the adjustment on
these characteristics with log assets and equity-to-assets included,
no specification explains even 1~percent of the variation; unrealized
securities losses carry a coefficient of 0.016, marginally significant
alone and insignificant once entered jointly with the brokered share
and CRE concentration; the uninsured deposit share enters
\emph{negatively} ($-0.003$, $p<0.05$), so if anything high-CECL banks
funded themselves with slightly fewer of the deposits that ran in
March 2023; and the composite itself enters negatively as well
($-0.406$, $p<0.01$).  A stress-exposure story for the results
requires the treatment to be correlated with stress exposure.  It is
not, and where a relation is detectable at all it runs the wrong way.

A placebo test turns the question around.  If depositors were reacting
to stress exposure rather than to the disclosure, uninsured time
deposits should have fallen at loss-exposed banks even where no CECL
news arrived.  Table~\ref{tab:svb_placebo} estimates a
difference-in-differences around the March 2023 failures in which the
treatment intensity is the 2022Q4 unrealized securities loss itself.
Columns (1)--(3) restrict the sample to the banks that reported a
Day-One adjustment of exactly zero, so the disclosure carried no news
about them.  Uninsured time deposits at loss-exposed banks did not
fall after the failures: the interaction is positive (0.073,
significant at 5~percent), and insured time deposits move the same way
(0.062), a parallel drift in the funding mix of banks holding
long-duration securities as rates rose, not a run.  Columns (4) and
(5) return to the full sample and add the loss interaction to the
baseline specification: the CECL coefficients ($-0.078$ continuous,
$-0.406$ High-CECL, both significant at 5~percent) are essentially
identical to the baseline estimates in
Table~\ref{tab:small_bank_did_time_deposits}, while the loss
interaction is small and insignificant.  The depositors who moved on
the disclosed charge did not move on the balance-sheet exposure that
defined the SVB episode.

The broader accounts of the episode point the same way.  The runs
were concentrated in crypto- and venture-focused business models at a
handful of mostly West Coast banks, and smaller banks faced less
pressure \citep{kelly2025rushing}; the deposit reallocation of the
stress weeks ran toward the largest banks on size and perceived
safety rather than on balance-sheet fundamentals
\citep{caglio2024depositors}; and the runs themselves were wholesale,
with no retail run and almost no run risk at private banks
\citep{cipriani2024tracing}.  None of these mechanisms operates
through a bank's CECL Day-One adjustment.
